\documentclass[10pt,aspectratio=169]{beamer}
% Preámbulo modular para presentaciones Beamer (Modelación Estadística)
% Diseño y paleta institucional del Tecnológico de Monterrey

\documentclass[aspectratio=169,xcolor={dvipsnames,table}]{beamer}

\usepackage[utf8]{inputenc}
\usepackage[spanish,mexico]{babel}
\usepackage{amsmath,amssymb,amsfonts,mathtools}
\usepackage{graphicx}
\usepackage{listings}
\usepackage{booktabs}
\usepackage{tikz}
\usetikzlibrary{matrix,arrows,decorations.pathmorphing,shapes.geometric,positioning}

% Paleta Institucional Tecnológico de Monterrey
\definecolor{TecAzul}{HTML}{0039A6}       % Azul TEC: color primario institucional
\definecolor{TecAzulOscuro}{HTML}{1A2E51} % Azul marino oscuro
\definecolor{TecRojo}{HTML}{EC2661}       % Rojo de acento (paleta miTec)
\definecolor{TecGris}{HTML}{646464}       % Gris neutro para comentarios y subtítulos
\definecolor{TecFondoBloque}{HTML}{F4F6F9}% Fondo suave para bloques

% Configuración de Tema Beamer
\usetheme{Boadilla}
\usecolortheme[named=TecAzulOscuro]{structure}
\setbeamercolor{alerted text}{fg=TecRojo}
\setbeamercolor{palette primary}{bg=TecAzulOscuro,fg=white}
\setbeamercolor{palette secondary}{bg=TecAzul,fg=white}
\setbeamercolor{palette tertiary}{bg=TecAzulOscuro!80!black,fg=white}
\setbeamercolor{title}{fg=TecAzulOscuro,bg=TecFondoBloque}
\setbeamercolor{frametitle}{fg=TecAzulOscuro,bg=TecFondoBloque}

% Bloques personalizados elegantes
\setbeamercolor{block title}{bg=TecAzulOscuro,fg=white}
\setbeamercolor{block body}{bg=TecFondoBloque,fg=black}
\setbeamercolor{block title alerted}{bg=TecRojo,fg=white}
\setbeamercolor{block body alerted}{bg=TecRojo!8,fg=black}
\setbeamercolor{block title example}{bg=TecAzul,fg=white}
\setbeamercolor{block body example}{bg=TecAzul!8,fg=black}

% Redondear bordes y añadir sombra sutil
\setbeamertemplate{blocks}[rounded][shadow=true]
\setbeamertemplate{navigation symbols}{}

% Pie de página limpio con número de diapositiva
\setbeamertemplate{footline}{
  \leavevmode%
  \hbox{%
  \begin{beamercolorbox}[wd=.7\paperwidth,ht=2.25ex,dp=1ex,left]{author in head/foot}%
    \hspace*{2ex}\insertshortauthor~~(\insertshortinstitute) \hfill \insertshorttitle
  \end{beamercolorbox}%
  \begin{beamercolorbox}[wd=.3\paperwidth,ht=2.25ex,dp=1ex,right]{date in head/foot}%
    \insertshortdate{}\hspace*{2em}
    \insertframenumber{} / \inserttotalframenumber\hspace*{2ex}
  \end{beamercolorbox}}%
  \vskip0pt%
}

% Configuración de Listings de Python para visualización pedagógica de código
\lstset{
    language=Python,
    basicstyle=\ttfamily\footnotesize,
    keywordstyle=\color{TecAzulOscuro}\bfseries,
    stringstyle=\color{TecRojo},
    commentstyle=\color{TecGris}\itshape,
    showstringspaces=false,
    breaklines=true,
    frame=lines,
    rulecolor=\color{TecAzulOscuro!30},
    backgroundcolor=\color{TecFondoBloque},
    aboveskip=0.5em,
    belowskip=0.5em,
    literate={á}{{\'a}}1 {é}{{\'e}}1 {í}{{\'i}}1 {ó}{{\'o}}1 {ú}{{\'u}}1
             {Á}{{\'A}}1 {É}{{\'E}}1 {Í}{{\'I}}1 {Ó}{{\'O}}1 {Ú}{{\'U}}1
             {ñ}{{\~n}}1 {Ñ}{{\~N}}1 {¿}{{?`}}1 {¡}{{!`}}1
}

% Metadatos institucionales por defecto
\title[Modelación Estadística]{Modelación Estadística para Ciencia de Datos e Ingeniería}
\author[J. Castillo Colmenares]{Juliho David Castillo Colmenares}
\institute[Tec de Monterrey]{Tecnológico de Monterrey \\ Escuela de Ingeniería y Ciencias}
\date{\today}

% Comandos y macros estadísticas compartidas para las presentaciones Beamer (Metropolis)

% Conjuntos numéricos básicos
\newcommand{\R}{\mathbb{R}}
\newcommand{\N}{\mathbb{N}}
\newcommand{\Q}{\mathbb{Q}}
\newcommand{\Z}{\mathbb{Z}}
\newcommand{\set}[1]{\left\{ #1 \right\}}
\newcommand{\abs}[1]{\left|#1\right|}
\newcommand{\norm}[1]{\left\| #1 \right\|}

% Comandos estadísticos y probabilísticos del curso
\newcommand{\Var}{\operatorname{Var}}
\newcommand{\cov}{\operatorname{Cov}}
\newcommand{\corr}{\rho}
\newcommand{\comb}[2]{\begin{pmatrix} #1 \\ #2 \end{pmatrix}}
\newcommand{\s}{\sigma}
\newcommand{\particion}{\mathfrak{P}}
\newcommand{\card}[1]{\#\left( #1 \right)}
\newcommand{\seq}[3]{\set{#1}_{#2}^{#3}}
\newcommand{\E}{\mathbb{E}}
\newcommand{\Prob}{\mathbb{P}}

% Entornos pedagógicos y cajas en español académico natural
\newenvironment{discusion}{%
  \begin{alertblock}{Pregunta de Discusión}%
}{%
  \end{alertblock}%
}

\newenvironment{reflexion}{%
  \begin{alertblock}{Reflexión para el Aula}%
}{%
  \end{alertblock}%
}

\newenvironment{paradoja}{%
  \begin{alertblock}{Paradoja de Exploración}%
}{%
  \end{alertblock}%
}

\newenvironment{motivacion}{%
  \begin{exampleblock}{Motivación: Ciencia de Datos en la Práctica}%
}{%
  \end{exampleblock}%
}

\newenvironment{retoclase}{%
  \begin{block}{Reto Rápido en Equipo}%
}{%
  \end{block}%
}


\title[Conceptos fundamentales]{Sección 5.8: Conceptos fundamentales de estadística inferencial}
\subtitle{Estimadores, error estándar y límites muestrales}
\author[J. Castillo Colmenares]{Juliho Castillo Colmenares}
\institute{Tecnológico de Monterrey}
\date{}

\begin{document}
\begin{frame}[plain]\vspace{-0.8cm}\titlepage\end{frame}

\begin{frame}{Hoja de ruta --- capítulo 5}
  \footnotesize
  \begin{itemize}\itemsep=0.08cm
    \item \textbf{05.00} Introducción inferencial
    \item \textbf{05.01--05.05} Distribuciones muestrales
    \item \textbf{05.08} Conceptos fundamentales \emph{(hoy)}
    \item \textbf{05.09} Estadísticos Z y t
  \end{itemize}
  \begin{block}{Objetivo didáctico}
    Relacionar estimadores, distribuciones muestrales y error estándar con la
    ley de los grandes números y el teorema del límite central.
  \end{block}
\end{frame}

\begin{frame}{Motivación: estimar un parámetro}
  \footnotesize
  Un estimador es un estadístico que se utiliza para aproximar un parámetro
  poblacional desconocido. Se escribe $\hat\theta$ para el parámetro $\theta$.
  \begin{alertblock}{Pregunta guía}
    ¿El estimador está centrado en el valor verdadero y qué tan variable es?
  \end{alertblock}
\end{frame}

\begin{frame}{Estimador insesgado}
  \footnotesize
  \begin{block}{Definición}
    $\hat\theta$ es insesgado para $\theta$ si
    \[
      E(\hat\theta)=\theta.
    \]
  \end{block}
  \begin{exampleblock}{Ejemplo de la nota}
    La media muestral cumple $E(\bar X)=\mu$.
  \end{exampleblock}
\end{frame}

\begin{frame}{Distribución muestral}
  \footnotesize
  La distribución muestral de un estadístico es su distribución de probabilidad
  al calcularlo sobre todas las muestras posibles de tamaño $n$.
  \begin{block}{Objeto de estudio}
    No describe una sola muestra: describe cómo cambia el estadístico al
    repetir el muestreo.
  \end{block}
\end{frame}

\begin{frame}{Error estándar}
  \footnotesize
  El error estándar es la desviación estándar de la distribución muestral.
  \[
    SE(\bar X)=\frac{\sigma}{\sqrt n},
    \qquad SE(\bar X)\approx\frac{s}{\sqrt n}\quad(\sigma\text{ desconocida}).
  \]
  \begin{alertblock}{Lectura}
    Mide la precisión de la estimación, no la dispersión de las observaciones
    individuales.
  \end{alertblock}
\end{frame}

\begin{frame}{Ley de los grandes números}
  \footnotesize
  Para variables iid con media $\mu$,
  \[
    \lim_{n\to\infty}P(|\bar X_n-\mu|<\varepsilon)=1.
  \]
  \begin{block}{Interpretación}
    Con muestras grandes, la media muestral se acerca en probabilidad a la
    media poblacional.
  \end{block}
\end{frame}

\begin{frame}{Teorema del límite central}
  \footnotesize
  Si $X_i$ son iid, tienen media $\mu$ y varianza finita $\sigma^2$, entonces
  \[
    \frac{\bar X-\mu}{\sigma/\sqrt n}\xrightarrow{d}N(0,1).
  \]
  En la práctica suele funcionar bien desde $n\ge30$, dependiendo de la forma
  de la población.
\end{frame}

\begin{frame}{Importancia para la inferencia}
  \footnotesize
  El TLC permite:
  \begin{itemize}\itemsep=0.1cm
    \item usar la normal para medias aun con poblaciones no normales;
    \item justificar pruebas Z y t;
    \item construir intervalos de confianza;
    \item comprender por qué aparecen formas aproximadamente normales.
  \end{itemize}
\end{frame}

\begin{frame}{Puente computacional 1/4: estimador}
  \footnotesize
  \[
    \bar X=\frac1n\sum_{i=1}^nX_i,
    \qquad E(\bar X)=\mu.
  \]
  \begin{block}{Verificación}
    El valor esperado identifica el centro del estimador con el parámetro que
    queremos estimar.
  \end{block}
\end{frame}

\begin{frame}{Puente computacional 2/4: variabilidad}
  \footnotesize
  \[
    \operatorname{Var}(\bar X)=\frac{\sigma^2}{n},
    \qquad SE(\bar X)=\frac{\sigma}{\sqrt n}.
  \]
  Aumentar $n$ reduce la variabilidad de la media muestral.
\end{frame}

\begin{frame}{Puente computacional 3/4: estandarización}
  \footnotesize
  \[
    Z_n=\frac{\bar X-\mu}{\sigma/\sqrt n}.
  \]
  \begin{alertblock}{Límite}
    El TLC afirma que $Z_n$ se aproxima a una normal estándar, lo que permite
    consultar probabilidades y cuantiles conocidos.
  \end{alertblock}
\end{frame}

\begin{frame}{Puente computacional 4/4: tamaño de muestra}
  \footnotesize
  \begin{tabular}{lll}
    \hline
    Cambio & Efecto en $SE(\bar X)$ & Lectura \\
    \hline
    $n\to4n$ & se divide entre $2$ & mayor precisión \\
    $n\to n/4$ & se multiplica por $2$ & menor precisión \\
    \hline
  \end{tabular}
  \begin{block}{Conclusión}
    La precisión mejora con $n$, aunque no elimina errores no muestrales.
  \end{block}
\end{frame}

\begin{frame}{Ejemplo resuelto 1/4 --- media insesgada}
  \footnotesize
  La nota considera $\bar X=(1/n)\sum X_i$ como estimador de $\mu$.
  \begin{block}{Planteamiento}
    Calcular $E(\bar X)$ usando linealidad de la esperanza.
  \end{block}
\end{frame}

\begin{frame}{Ejemplo resuelto 1/4 --- solución}
  \footnotesize
  \[
    E(\bar X)=E\!\left(\frac1n\sum X_i\right)
    =\frac1n\sum E(X_i)=\frac1n\,n\mu=\mu.
  \]
  \begin{alertblock}{Resultado}
    La media muestral es un estimador insesgado de la media poblacional.
  \end{alertblock}
\end{frame}

\begin{frame}{Ejemplo resuelto 2/4 --- distribución muestral}
  \footnotesize
  Para variables iid con media $\mu$ y varianza $\sigma^2$, la nota estudia la
  distribución de $\bar X$.
  \begin{block}{Planteamiento}
    Determinar su media y su varianza.
  \end{block}
\end{frame}

\begin{frame}{Ejemplo resuelto 2/4 --- solución}
  \footnotesize
  \[
    E(\bar X)=\mu,\qquad \operatorname{Var}(\bar X)=\frac{\sigma^2}{n}.
  \]
  \begin{exampleblock}{Interpretación}
    La variabilidad de la media disminuye a medida que aumenta el tamaño de la
    muestra.
  \end{exampleblock}
\end{frame}

\begin{frame}{Ejemplo resuelto 3/4 --- tiempos exponenciales}
  \footnotesize
  El tiempo de espera tiene distribución exponencial con media $\mu=5$ minutos.
  Para $n=50$, la nota aplica el TLC a la media muestral.
  \begin{block}{Planteamiento}
    Encontrar la distribución aproximada de $\bar X$.
  \end{block}
\end{frame}

\begin{frame}{Ejemplo resuelto 3/4 --- solución}
  \footnotesize
  \[
    \bar X\approx N\!\left(5,\frac{5^2}{50}\right)=N(5,0.5).
  \]
  \begin{alertblock}{Interpretación}
    Aunque la población es exponencial, la media de 50 observaciones se modela
    aproximadamente con una normal.
  \end{alertblock}
\end{frame}

\begin{frame}{Aplicación de lectura 4/4 --- solución}
  \footnotesize
  \[
    \text{muestra}\to\text{estadístico}\to\text{distribución muestral}
    \to SE\to\text{aproximación normal}\to\text{inferencia}.
  \]
  Esta secuencia reproduce la estructura conceptual de la nota.
\end{frame}

\begin{frame}{Síntesis de la sección}
  \footnotesize
  \begin{itemize}\itemsep=0.12cm
    \item La insesgadez centra el estimador.
    \item El error estándar cuantifica su precisión.
    \item La ley de los grandes números explica la convergencia.
    \item El TLC permite trabajar con una normal límite.
  \end{itemize}
\end{frame}

\begin{frame}{Transición curricular}
  \footnotesize
  \begin{block}{Lo que queda establecido}
    Ya podemos describir cómo se comporta un estimador al repetir el muestreo.
  \end{block}
  \begin{alertblock}{Siguiente sección}
    \textbf{05.09 Estadísticos Z y t}: estandarización para pruebas sobre medias
    con desviación conocida o desconocida.
  \end{alertblock}
\end{frame}
\end{document}
